% Part: first-order-logic
% Chapter: model-theory
% Section: partial-iso

\documentclass[../../../include/open-logic-section]{subfiles}

\begin{document}

\olfileid{mod}{bas}{pis}
\section{Een isomorfisme dat maar op een eindig deel is ingevuld}

\begin{defn}
  Neem twee !!{structure}s $\Struct{M}$ en $\Struct{N}$. Een
  \emph{partieel isomorfisme} van $\Struct{M}$ naar $\Struct{N}$ is
  een functie $p$ die maar voor eindig veel elementen van $\Domain M$
  is ingevuld en die waarden in $\Domain N$ geeft. Op dat eindige deel
  moet de functie zich houden aan de voorwaarden voor een isomorfisme uit
  \olref[iso]{defn:isomorphism}:
  \begin{enumerate}
  \item $p$ is !!{injective}: twee verschillende invoeren krijgen dus
    nooit dezelfde uitkomst;
  \item neem een !!{constant}~$c$. Als $p(\Assign{c}{M})$ is
    gedefinieerd, dan $p(\Assign{c}{M}) = \Assign{c}{N}$;
  \item neem een $n$-plaatsig !!{predicate} $P$. Als $a_1$, \dots, $a_n$
    in het domein van $p$ zitten, dan geldt
    $\langle a_1, \dots, a_n\rangle \in \Assign P M$ precies wanneer
    $\langle p(a_1), \dots, p(a_n) \rangle \in \Assign P N$;
  \item neem een $n$-plaatsige !!{function} $f$. Als $a_1$, \dots, $a_n$
    in het domein van $p$ zitten, dan $p(\Assign f M (a_1, \dots,a_n))
    = \Assign f N (p(a_1), \dots, p(a_n))$.
  \end{enumerate}
  Met ‘$p$ is eindig’ bedoelen we dus dat $\dom{p}$ eindig is.
\end{defn}

De lege functie~$\emptyset$ is altijd een partieel isomorfisme tussen
om het even welke twee !!{structure}s. Er hoeft dan immers nog geen
enkel element te worden vergeleken.

\begin{defn}\ollabel{defn:partialisom}
  Twee !!{structure}s $\Struct{M}$ en $\Struct{N}$ zijn
  \emph{partieel isomorf}, geschreven als $\Struct{M} \iso[p]
  \Struct{N}$, als er een niet-lege verzameling $I$ van partiële
  isomorfismen tussen $\Struct{M}$ en $\Struct{N}$ is waarmee je aan
  beide kanten steeds verder kunt gaan. Precies gezegd moet die verzameling aan de
  \emph{heen-en-weereigenschap} voldoen:
  \begin{enumerate}
  \item (\emph{Heen}) Neem een $p \in I$ en een $a \in \Domain M$.
    Dan is er een $q \in I$ met $p \subseteq q$ waarbij $a$ in het
    domein van $q$ zit;
  \item (\emph{Weer}) Neem een $p \in I$ en een $b \in \Domain N$.
    Dan is er een $q \in I$ met $p \subseteq q$ waarbij $b$ in het
    bereik van $q$ zit.
  \end{enumerate}
\end{defn}

\begin{thm}\ollabel{thm:p-isom1}
  Als $\Struct{M} \iso[p] \Struct{N}$ en $\Struct{M}$ en
  $\Struct{N}$ !!{enumerable} zijn, dan $\Struct{M} \iso
  \Struct{N}$.
\end{thm}

\begin{proof}
  Omdat $\Struct{M}$ en $\Struct{N}$ !!{enumerable} zijn, kunnen we hun
  elementen in rijtjes zetten:
  $\Domain{M} = \{a_0, a_1, \ldots \}$ en
  $\Domain{N} = \{b_0, b_1, \ldots \}$. Kies eerst zomaar een
  $p_0 \in I$. Daarna maken we stap voor stap een steeds grotere rij
  partiële isomorfismen
  $p_0 \subseteq p_1 \subseteq p_2 \subseteq \cdots$:
  \begin{enumerate}
  \item als $n+1$ oneven is, dus $n = 2r$, gebruiken we de
    heen-regel. Daarmee kiezen we $p_{n+1} \in I$ zo dat
    $p_n \subseteq p_{n+1}$ en $a_r$ in het domein van $p_{n+1}$ zit;
  \item als $n+1$ even is, dus $n+1 =2r$, gebruiken we de
    weer-regel. Daarmee kiezen we $p_{n+1} \in I$ zo dat
    $p_n \subseteq p_{n+1}$ en $b_r$ in het bereik van $p_{n+1}$ zit.
  \end{enumerate}
Voeg nu alle stappen samen:
\[
p = \bigcup_{n\ge 0} p_n,
\]
De oneven stappen zorgen dat elk element van de eerste structuur uiteindelijk
in het domein komt. De even stappen zorgen dat elk element van de tweede
structuur uiteindelijk in het bereik komt. Zo is $p$ een
isomorfisme tussen $\Struct{M}$ en $\Struct{N}$.
\end{proof}

\begin{prob}
  Werk in detail uit waarom de $p$ uit
  \olref[mod][bas][pis]{thm:p-isom1} inderdaad een isomorfisme is.
\end{prob}

\begin{thm}\ollabel{thm:p-isom2}
  Stel dat $\Struct{M}$ en $\Struct{N}$ !!{structure}s zijn voor een
  zuiver relationele !!{language}. Dat is hier !!a{language} met
  alleen !!{predicate}s en zonder !!{function}s of constanten. Als
  $\Struct{M} \iso[p] \Struct{N}$, dan geldt ook
  $\Struct{M} \elemequiv \Struct{N}$.
\end{thm}

\begin{proof}
  We bewijzen dit met inductie over !!{formula}s. Neem elementen
  $a_1$, \dots, $a_n$ en $b_1$, \dots, $b_n$. Stel dat een partieel
  isomorfisme $p$ elk $a_i$ op $b_i$ afbeeldt. Neem verder bedelingen
  met $s_1(x_i) =a_i$ en $s_2(x_i) =b_i$ (voor
  $i =1$, \dots,~$n$). De inductie laat dan zien dat
  $\Sat{M}{!A}[s_1]$ precies wanneer $\Sat{N}{!A}[s_2]$. Bij $n=0$
  hebben we geen gekozen elementen meer. Dan volgt
  $\Struct{M} \elemequiv \Struct{N}$.
\end{proof}

\begin{rem}
Als de taal wel !!{function}s heeft, blijft het vorige resultaat waar.
Dan moet je kijken naar het isomorfisme dat $p$ geeft tussen de
sub!!{structure} van $\Struct{M}$ die door $a_1$, \dots, $a_n$ wordt
voortgebracht en de sub!!{structure} van $\Struct{N}$ die door
$b_1$, \dots, $b_n$ wordt voortgebracht.
\end{rem}

We kunnen het vorige resultaat in stappen opdelen. Het aantal kwantoren
dat in !!a{formula} in elkaar staat, bepaalt hoe vaak we de bijbehorende
partiële isomorfismen moeten kunnen uitbreiden.

\begin{defn}
  Neem een !!{formula}~$!A$. De \emph{kwantordiepte} van $!A$,
  geschreven als $\QuantRank{!A} \in \Nat$, is het grootste aantal
  kwantoren dat in $!A$ in elkaar staat. We bepalen dat aantal stap
  voor stap uit de opbouw van de formule. Twee !!{structure}s
  $\Struct{M}$ en $\Struct{N}$ zijn \emph{$n$-equivalent}, geschreven
  als $\Struct{M} \elemequiv[n] \Struct{N}$, als zij dezelfde zinnen
  waar maken zolang de kwantordiepte kleiner dan of gelijk aan~$n$ is.
\end{defn}

\begin{prop}\ollabel{prop:qr-finite}
  Neem een eindige zuiver relationele !!{language} $\Lang{L}$. Daarmee
  bedoelen we hier !!a{language} met eindig veel !!{predicate}s en
  !!{constant}s, maar zonder !!{function}s. Voor elke $n \in \Nat$
  zijn er, als logisch equivalente zinnen als hetzelfde tellen, maar
  eindig veel eerste-orde!!{sentence}s in de !!{language} $\Lang{L}$
  waarvan de kwantordiepte op zijn meest $n$ is.
\end{prop}

\begin{proof}
  Met inductie over $n$.
\end{proof}

\begin{defn}
  Neem !!a{structure}~$\Struct{M}$. De verzameling
  $\Domain M^{<\omega}$ bestaat uit alle eindige rijtjes van elementen
  van $\Domain{M}$. Met $\mathbf{a}, \mathbf{b}, \mathbf{c}, \ldots$
  bedoelen we zulke rijtjes. Als
  $\mathbf{a} \in \Domain{M}^{<\omega}$ en $a \in \Domain{M}$, dan
  betekent $\mathbf{a}a$ de \emph{concatenatie} van $\mathbf{a}$ met
  $a$: we zetten het laatste element dus achter het rijtje.
\end{defn}

\begin{defn}
  Neem !!{structure}s $\Struct{M}$ en $\Struct{N}$. We maken relaties
  $I_n \subseteq \Domain M^{<\omega} \times \Domain N^{<\omega}$
  tussen rijtjes die even lang zijn. Voor elke $n$ doen we dat als volgt:
   \begin{enumerate}
   \item $I_0(\mathbf{a},\mathbf{b})$ geldt precies wanneer
     $\mathbf{a}$ en $\mathbf{b}$ in $\Struct{M}$ en $\Struct{N}$
     aan dezelfde atomaire !!{formula}s voldoen. Anders gezegd: stel
     $s_1(x_i) = a_i$ en $s_2(x_i) = b_i$. Als $!A$ atomair is en alle
     !!{variable}s ervan bij $x_1$, \dots,~$x_n$ staan, dan geldt
     $\Sat{M}{!A}[s_1]$ precies wanneer~$\Sat{N}{!A}[s_2]$;
   \item $I_{n+1} (\mathbf{a},\mathbf{b})$ geldt precies wanneer je
     voor elke $a\in \Domain M$ een $b\in \Domain N$ kunt vinden met
     $I_n(\mathbf{a}a,\mathbf{b}b)$, en ook omgekeerd.
   \end{enumerate}
\end{defn}


\begin{defn}
  We schrijven $\Struct{M} \approx_n \Struct{N}$ als
  $I_n(\emptyseq,\emptyseq)$ geldt voor $\Struct{M}$ en
  $\Struct{N}$. Hier is $\emptyseq$ het lege rijtje.
\end{defn}

\begin{thm}\ollabel{thm:b-n-f}
  Neem een zuiver relationele !!{language} $\Lang{L}$. Als
  $I_n(\mathbf{a},\mathbf{b})$, dan geldt voor elke $!A$ met
  $\QuantRank{!A} \le n$ dat $\Sat{M}{!A}[\mathbf{a}]$ precies wanneer
  $\Sat{N}{!A}[\mathbf{b}]$. Hier zeggen we weer dat $\mathbf{a}$ aan
  $!A$ voldoet als elke bedeling $s$ met $s(x_i) = a_i$ aan $!A$
  voldoet. Als $\Lang{L}$ eindig is, geldt de omgekeerde richting ook.
\end{thm}

\begin{proof}
  Uit $I_n(\mathbf{a},\mathbf{b})$ volgt dat $\mathbf{a}$ en
  $\mathbf{b}$ aan dezelfde !!{formula}s met kwantordiepte op zijn
  meest $n$ voldoen. Dat bewijzen we met een gewone inductie over $!A$.
  Voor de omgekeerde richting gebruiken we inductie over $n$. Daarbij
  gebruiken we \olref{prop:qr-finite}: voor elke $n$ zijn er maar
  eindig veel !!{formula}s van die kwantordiepte die niet logisch
  equivalent zijn.

  Neem eerst $n=0$. Als $\mathbf{a}$ en $\mathbf{b}$ aan dezelfde
  kwantorvrije !!{formula}s voldoen, voldoen zij ook aan dezelfde
  atomaire formules. Dus $I_0(\mathbf{a},\mathbf{b})$.

  Neem nu het geval $n+1$. Stel dat $\mathbf{a}$ en $\mathbf{b}$ aan
  dezelfde !!{formula}s met kwantordiepte op zijn meest $n+1$ voldoen.
  Voor $I_{n+1}(\mathbf{a},\mathbf{b})$ moeten we het volgende laten
  zien: bij elke $a \in \Domain M$ is een $b \in \Domain N$ te vinden
  waarvoor $I_n(\mathbf{a}a,\mathbf{b}b)$. Door de inductiehypothese is
  het weer genoeg te laten zien dat er voor elke $a \in \Domain M$ een
  $b \in \Domain N$ is waarvoor $\mathbf{a}a$ en $\mathbf{b}b$ aan
  dezelfde !!{formula}s met kwantordiepte op zijn meest $n$ voldoen.

  Kies $a \in \Domain M$. Laat $!T^a_n$ de verzameling zijn van alle
  !!{formula}s $!B(x,\mathbf{y})$ met kwantordiepte op zijn meest $n$
  waaraan $\mathbf{a}a$ in $\Struct{M}$ voldoet. Tot op logische
  equivalentie zijn dat er maar eindig veel. Neem van die eindige
  verzameling de conjunctie; zo kunnen we $!T^a_n$ als één
  eerste-orde!!{formula} behandelen. De rij $\mathbf{a}$ voldoet dan
  aan $\lexists[x][!T^a_n(x,\mathbf{y})]$. Deze formule heeft
  kwantordiepte op zijn meest $n+1$. Volgens onze aanname voldoet
  $\mathbf{b}$ in $\Struct{N}$ aan dezelfde !!{formula}. Er is dus een
  $b \in \Domain N$ waarvoor $\mathbf{b}b$ aan $!T^a_n$ voldoet.
  Daarmee voldoet $\mathbf{b}b$ aan dezelfde !!{formula}s met
  kwantordiepte op zijn meest $n$ als $\mathbf{a}a$.
  Draai nu de rollen om. Dan vinden we voor elke $b \in \Domain N$ een
  $a\in \Domain M$ waarvoor $\mathbf{a}a$ en $\mathbf{b}b$ aan
  dezelfde !!{formula}s met kwantordiepte op zijn meest $n$ voldoen.
  Dat zijn precies de twee richtingen die we nodig hadden.
\end{proof}

\begin{cor}\ollabel{cor:b-n-f}
  Als $\Struct{M}$ en $\Struct{N}$ zuiver relationele !!{structure}s
  in een eindige !!{language} zijn, dan geldt
  $\Struct{M} \approx_n\Struct{N}$ precies wanneer
  $\Struct{M} \elemequiv[n] \Struct{N}$. In het bijzonder geldt
  $\Struct{M} \elemequiv \Struct{N}$ precies wanneer voor elke $n$
  geldt dat $\Struct{M} \approx_n \Struct{N}$.
\end{cor}

\end{document}
